% =====================================================================
%  APPENDICES (Annexure 12): labelled "APPENDIX A: TITLE"
% =====================================================================
\clearpage
\setcounter{chapter}{0}
\renewcommand{\thechapter}{\Alph{chapter}}
\chapter*{APPENDIX A: CONTRIBUTION OF PROJECT}
\refstepcounter{chapter}
\phantomsection
\addcontentsline{toc}{chapter}{APPENDIX A: CONTRIBUTION OF PROJECT}

\section*{A.1 Objective and Relevance of Project}

The objective of \veriscript{} is to give institutions an open, explainable and
locally run tool that estimates whether a piece of English prose was written by
a person or generated by a large language model, and that makes the cost,
accuracy and fairness of that estimate visible. Large language models have made
it possible to produce fluent essays, reports, reviews and answers in seconds.
Teachers, examiners, admissions officers, editors and moderators now need to
decide whether submitted text reflects the effort of its stated author, but the
tools available to them are either commercial black boxes that charge per page
and do not explain their decisions, or research prototypes that are accurate on
one benchmark and fragile elsewhere.

The project addresses this gap with a two-stage design. A heuristic stage
computes interpretable statistics (how predictable the text is to a reference
language model, how uniform its sentences are, how varied its vocabulary is)
and decides confidently on most documents, explaining each decision with the
features that drove it. Only the documents on which it is uncertain are passed
to a fine-tuned transformer. The project also contributes an evaluation that
goes beyond a single accuracy figure: it measures the accuracy of every system
with confidence intervals, its behaviour on an unseen model family, on
out-of-distribution text and on adversarial edits, its false positive rate on
native and non-native student essays, and its latency on CPU and GPU.

The relevance of the work is practical. A false accusation of using AI can harm
a student's academic record; the base-rate argument of Chapter~1 shows that even
a detector that is ``95\% accurate'' produces many such accusations when AI use
is uncommon. By reporting probabilities, explanations and false positive rates,
and by showing where a fine-tuned detector breaks down, \veriscript{} helps an
institution use detection as a prompt for human review rather than as a verdict.

\section*{A.2 Technical Novelty and Utility}

The technical novelty of the project lies in four points.

\textbf{A confidence-gated cascade for AI-text detection.} Cascades and early
exit are established in computer vision and in language-model serving, but the
surveyed detection systems use either statistical features or a fine-tuned
encoder, not both in a calibrated routing scheme. \veriscript{} formulates the
choice of the ambiguity band as a constrained optimisation (minimum escalation
subject to an accuracy tolerance and a fairness constraint), solves it exactly
by a vectorised search over 2,601 candidate bands on cached predictions, and
exposes the result as a menu of operating points.

\textbf{A fairness-aware tuning procedure.} The project shows empirically that a
constraint on the disparity of false positive rates can be satisfied in a
degenerate way, by making a detector equally bad for every group, and replaces
it with a worst-group constraint that cannot be gamed.

\textbf{A unified evaluation framework.} The same code path serves the
heuristic-only, transformer-only and cascade systems, so the comparison is exact.
The framework covers bootstrap confidence intervals and paired tests,
leave-one-generator-out retraining, out-of-distribution and adversarial sets,
fairness, calibration and per-stage latency, and regenerates every table and
figure of Chapter~7 from the command line.

\textbf{Measured negative results.} The project quantifies, component by
component, why a GPT-2-based heuristic stage is not cheaper than a RoBERTa
classifier, and that a fine-tuned transformer with near-perfect in-distribution
AUROC flags most student essays as AI-generated. Both findings are useful to
anyone designing detectors.

The utility of the system is that it can be used as it is: as a Python library,
a command line tool for batch screening, a REST service for integration with
other systems, and a web page for individual reviewers, all running locally,
with commands to retrain and recalibrate it on an institution's own data.

\section*{A.3 Deliverables}

The project delivers the following items, all contained in the project
repository:
\begin{itemize}
  \item the \texttt{veriscript} Python package (pre-processing, feature
        extraction, heuristic and transformer stages, confidence gate, cascade,
        explanations and evaluation toolkit);
  \item the trained models: a calibrated XGBoost heuristic head, a fine-tuned
        RoBERTa-base classifier, the tuned operating points, and the
        leave-one-generator-out variants;
  \item a reproducible data pipeline that downloads HC3, M4, RAID, MAGE and the
        fairness essays and builds the 44,618-document corpus, plus a script for
        topic-matched custom generations;
  \item scripts that train, tune and evaluate the systems and produce every
        result table and figure;
  \item a command line tool, a FastAPI REST service, a Streamlit web interface
        and a Dockerfile;
  \item a test suite of 34 automated tests;
  \item documentation: a technical README, a plain-language user manual, and
        this report with its LaTeX sources.
\end{itemize}

\section*{A.4 Concerns Related to the Project}

\subsection*{A.4a Social Relevance}

Written assessment is one of the main ways in which schools and universities
judge learning, and recruitment, publishing and online platforms depend on
trusting that text was written by the person who submits it. Generative AI has
weakened that trust for everyone, including the many students who do their own
work. A transparent detector can help restore it, but only if it is used
carefully: the results of this project show that detectors flag the writing of
non-native English speakers more often than that of native speakers, and that a
fine-tuned model can flag most essays of a genre it has not seen. For this
reason \veriscript{} reports a probability with its reasons, states on every
result that it must be reviewed by a person, and publishes its false positive
rates by group. Used in this way, it can direct a teacher's attention to
submissions that deserve a conversation, while protecting students from
automated accusations. Because it is open source and runs on ordinary hardware,
it is also available to institutions that cannot afford commercial licences.

\subsection*{A.4b Environmental Sustainability}

Every detector consumes energy each time it is run, and institutional screening
can involve tens of thousands of documents per term. The project therefore
measured cost explicitly and runs entirely on a single workstation; the full
pipeline, including fine-tuning RoBERTa-base twice, took a few hours of one
consumer GPU, and no cloud resources were used. The measurements also show that
the current heuristic stage does not save energy compared with the transformer,
and Chapter~8 identifies how a smaller reference model would make the cascade
cheaper than either stage alone. Models are trained once and reused; inference
runs on existing hardware, a CPU being sufficient; and the data pipeline
samples only the part of the large RAID benchmark that is needed. The project
thus treats computational efficiency as a first-class objective, in line with
efforts to reduce the energy footprint of machine learning.

\subsection*{A.4c Ethical, Legal and Cultural Aspects}

\textbf{Ethical.} The output of an AI-text detector can affect a person's grades,
admission or reputation. \veriscript{} is designed as decision support: it
returns calibrated probabilities and explanations, never a bare verdict, and it
documents its error rates, including the higher false positive rate for
non-native writers. It should not be used as the sole basis for any sanction.

\textbf{Legal and privacy.} Student work is personal data. The system processes
documents locally and stores nothing unless logging is explicitly enabled, so
no submission is sent to a third party. All datasets used are public research
benchmarks used under their licences, and all software and model weights are
openly licensed.

\textbf{Cultural.} Writing style varies with language background, education and
genre. A detector trained mostly on English web text and ChatGPT answers
encodes a particular notion of ``typical'' human writing, which disadvantages
writers whose style differs from it. The fairness evaluation makes this visible,
and the recalibration commands allow an institution to adapt the system to its
own population of writers.

% ---------------------------------------------------------------------
\clearpage
\chapter*{APPENDIX B: INSTALLATION AND USER GUIDE}
\refstepcounter{chapter}
\phantomsection
\addcontentsline{toc}{chapter}{APPENDIX B: INSTALLATION AND USER GUIDE}

\section*{B.1 Installation}

\veriscript{} requires Python 3.11; a CUDA-capable GPU is optional. The commands
below install the package and its dependencies in a virtual environment and run
the test suite (PowerShell on Windows; the same commands work in a Unix shell
with \texttt{.venv/bin/python}).

\begin{lstlisting}[language={},numbers=none,caption={Installation},label={lst:install}]
py -3.11 -m venv .venv
.venv\Scripts\python -m pip install torch==2.5.1 `
    --index-url https://download.pytorch.org/whl/cu124   # or /whl/cpu
.venv\Scripts\python -m pip install -e ".[dev]"
.venv\Scripts\python -m spacy download en_core_web_sm
.venv\Scripts\python -m pytest -q
\end{lstlisting}

The trained models are produced by the pipeline of Listing~\ref{lst:pipeline}
and stored in the \texttt{models/} folder.

\section*{B.2 Checking Documents}

\begin{lstlisting}[language={},numbers=none,caption={Checking documents from the command line},label={lst:userguide}]
veriscript detect essay.docx                    # one file
veriscript detect --text "..." --mode thorough  # pasted text
veriscript batch submissions/ --out report.csv  # a folder or CSV
veriscript demo                                  # web page
veriscript serve                                 # REST API on :8000
\end{lstlisting}

The \texttt{--mode} option selects the operating point: \texttt{fast} and
\texttt{balanced} decide almost every document with the heuristic stage, and
\texttt{thorough} escalates about 8\% of documents and gives the most accurate
results. The batch report contains, for every document, the label, the
probability, the route, the latency, the top features and any note such as
``insufficient text''; a summary file gives the counts, the escalation rate and
the throughput.

\section*{B.3 Adapting the System to an Institution}

An administrator with a labelled set of the institution's own documents can
retrain and recalibrate the heuristic head and re-tune the operating points:

\begin{lstlisting}[language={},numbers=none,caption={Retraining and re-tuning on local data},label={lst:adapt}]
veriscript retrain-heuristic \
    --data local_labelled.csv \
    --out models/heuristic/local.joblib
veriscript tune --data local_validation.csv
\end{lstlisting}

The REST endpoint \texttt{PUT /api/v1/config} changes the operating point of a
running service. It requires the administrator token, which is set in the
environment variable \texttt{VERISCRIPT\_ADMIN\_TOKEN}.

% ---------------------------------------------------------------------
\clearpage
\chapter*{APPENDIX C: CONFIGURATION FILE}
\refstepcounter{chapter}
\phantomsection
\addcontentsline{toc}{chapter}{APPENDIX C: CONFIGURATION FILE}

All run-time settings are read from \texttt{configs/default.yaml}
(Listing~\ref{lst:config}). The operating points tuned by the threshold-tuning
script are stored in the file \url{models/operating_points.json} and override
the fall-back values in the configuration file.

\begin{lstlisting}[language={},numbers=none,caption={Default configuration (\texttt{configs/default.yaml})},label={lst:config}]
device: auto                 # auto | cpu | cuda
seed: 42
preprocess:
  min_words: 100             # shorter documents: "insufficient text"
reference_lm:
  name: gpt2                 # GPT-2 small, ~124M parameters
  max_tokens: 2048
  min_sentence_tokens: 5
heuristic:
  model_path: models/heuristic/heuristic.joblib
transformer:
  model_dir: models/transformer/roberta-base
  max_len: 512
  stride: 384                # step between window starts (tokens)
  batch_size: 16
gate:
  tau: 0.5                   # final decision threshold
  operating_point: balanced  # fast | balanced | thorough | custom
operating_points_path: models/operating_points.json
explain:
  top_k: 5
  note: "Probabilistic estimate; review before acting on it."
logging:
  enabled: false             # documents are not stored by default
\end{lstlisting}
